\section{Prerequisite NLP Information}
\label{sec:literature_prerequisite_nlp_information}
This section briefly defines and explains some of the fundamentals in NLP, which are referenced in the NLP-research section of this report's corresponding final-year- project document.
%TODO 89318
\subsection{Tokenisation}
Tokens are produced as a result of tokenisation, which is a process that breaks the input data -- text, for NLP -- into smaller units: tokens. There are different methods of tokenisation, for example breaking a sequence of text, e.g. a sentence, down to each word in the sentence (word tokenisation), an example of an array storing word-tokens is [``the'', ``fox'', ``jumped'', ``over'', ``the'', ``lazy'', ``dog'']. However, the method of tokenisation depends on the context and problem statement, for example it may be more appropriate to tokenise by subword, or characters in a (sub)word \citep{geeks:tokenisation}. Note, however, that whilst here and throughout this document, tokenisation examples are shown as they relate/apply to NLP, however tokenisation for computer vision, for example, may be a frame being chunked, each smaller chunk then being a token.
Tokenisation is used in NLP, since it allows for example a full-sentence to be broken down to individual components (words), which allows each word to be processed separately (in word-tokenisation), allowing for lexical analysis, such as the meaning of each word in context, NER, sentiment analysis, among others \citep{wikipedia:NLP}, examples of sentiment analyses for different NLP techniques are given in the FYP report, for example BoW in section \verb*|BoW_example| of FYP. %TODO more here pls
\subsection{Embedding}
Embedding is the process by which tokens are mapped to points in $N$-dimensional space, the position then stored in a vector, per token.
There are two types of resultant vectors produced by an embedding process (of which there are many, some of which are covered in \verb*|literature_existing_nlp_techniques| of the FYP main report) sparse embeddings, and dense embeddings \citep{zilliz:sparse_dense_embeddings}.
\subsubsection{Sparse-Embedding Matrix}
A sparse-embedding matrix is a matrix containing sparsely-embedded vectors for a given corpus, vectors embedded using a sparse method, such as one-hot encoding (such as BoW uses; see \verb*|para:bag_of_words| of the FYP main report). When using a sparse-embedding method, such as one-hot encoding, the number of dimensions will be the number of unique tokens (words, for NLP) in the corpus. With reference to the example showed in the BoW explanation in the FYP main report, in a given token's sparsely-embedded vector the majority of elements are zero, with a ``1'' element denoting the link between this unique-token and its corresponding position in the built lexicon, e.g. for a sparse-vector with contents [0,0,1] representing the token ``cat'' , ``cat'''s index in the arrayed-lexicon is 2 (zero-indexed). Since the number of dimensions are equal to the number of unique tokens, sparse embeddings are inherently very-high dimensionality, resulting in issues explained in \cref{subsubsec:curse_of_dimensionality}.
\paragraph{Use Cases}
\noindent\\
The structure of sparse-embeddings (large number of zeroed values) can be useful for keyword-lookup or matching, since the zero-values indicate an absence in their corresponding tokens, given the embedded-input \citep{zilliz:sparse_dense_embeddings}, which can be coupled with vector indexing to map vectors to their nearest-neighbours. This enables keyword-lookup for an inputted keyword, allowing similar present keywords [in the input-corpus] to be returned, resulting in fast similarity lookups \citep{zilliz:vector_indexing}, but only when a method of embedding involves associating a token with its nearest-neighbour, such as SPLADE.
\subparagraph{SPLADE}
\noindent\\
SPLADE uses modern NLP techniques, such as transformers (see \verb*|subsubsec:transformers| of the FYP main report), to generate a sparse-embedding vector of the input, where the transformer aspect captures semantically-similarities between the input-tokens (via self-attention: \verb*|lit:self-attention|), this can then be coupled with inverted indexing to map a specific keyword (-to-be looked-up) to corpora containing that specific keyword, or contextually-similar keywords, which may, for example, associate ``flights'' with ``airlines'', both pertaining to air-travel \citep{pinecone:splade}, this can hence be effective when used in search engines' back-end, for example, or other information-retrieval systems.
\subsubsection{Dense-Embedding Matrix}
Matrix-collation of densely-embedded vectors. Densely-embedded vectors are, hence the name, more compact than sparsely-embedded vectors, containing few (or no) zeroed-values. Crucially, however, the embedding process typically involves a neural-network approach (e.g. Word2Vec; BERT (see FYP main report corresponding sections for each Word2Vec and BERT)), allowing semantic-relationships to be captured in the resultant, ordered, embedding. Given the few-, or non-, zeroed-values nature of \textbf{dense} embeddings, these are inherently lower in dimensionality than sparse embeddings. The issues arising from the curse of dimensionality (\cref{subsubsec:curse_of_dimensionality}) do not apply here, given the lower dimensionality, as such they are faster to compute than sparse embeddings\footnote{It is worth noting that processing sparse embeddings tends to result in a RAM-bottleneck, as a result of the very-large number of (zeroed) elements; this is not the case for an equivalent-input dense embedding, as such these [dense embeddings] require less RAM bandwidth, compared with processing sparse embeddings, though dense-embeddings can be (and, actually, in practice, often are much) more computationally-expensive to map/create initially.}.\\
In NLP, dense embeddings are used to (significantly-better) capture semantic-meaning from the input. Therefore, unlike sparse-embeddings, these can be used for strong(er) text analysis, for example the sentiment of a review (customer's perceived satisfaction) can be contained within its dense embedding (unlike a bag of words approach with a sparse-embedding, particularly given BoW models/sparse-embedding vectors are unordered),_ this then an input to a downstream task, such as classification, wherein the review would be classified into categories, e.g. binary classification of good/bad review. With this, the dense embedding is learnt through a technique such as masked-language modelling (explained in the FYP main report).\\
This project will use dense embeddings, since the nature of NLP analysis required is such that the (potentially ambiguous) meaning (/intention) of a received communication can be understood, thereby allowing this sentiment to be classified into threat categories.

%TODO \textcolor{red}{REF and continue readeing this page: https://zilliz.com/learn/sparse-and-dense-embeddings}
\subsection{Dimensionality}
\label{subsec:dimensionality}
Dimensionality, informally, refers to the number of coordinates to define a point in a given-$N$- dimensional space \citep{wikipedia:dimensionality}. With embedding vectors, this is the length of the vector (number of elements).\\
Visualising embeddings in a two/three-dimensional space can be achieved by using a dimensionality-reduction technique such as Principal-Component Analysis (PCA), retaining the most important data-points, whilst reducing the number of features (aspects) of the embedding (see below \cref{para:reducing_dimensionality_via_pca}).

\subsubsection{Curse of Dimensionality}%TODO incomplete
\label{subsubsec:curse_of_dimensionality}
%having high num of dimensions could indicate sparsity (eg as in sparse-embeddings) from ig not being able to fill those dims. with data?
The curse of dimensionality, introduced by Bellman, refers to problems that arise in data with high dimensionality, for example a sharp decrease in the accuracy of a model when increasing the number of dimensions in the data, which also causes problems for data-acquisition itself, since the amount of data required (feature space) increases exponentially with increases in dimensionality. Further, computational expense of an increased amount of data, from an increase in features (potentially higher accuracy, however -- when building-up lower-dimensional data -- on downstream tasks).\\
Given that dimensionality is the amount of features/attributes comprising some data, e.g. aforementioned embedding vector, if the vector has particularly-high dimensionality, per additional vector computed, there are a particularly-high number of fields to fill, i.e. if all vectors have 250 dimensions, each vector must have 250 data points (entries) (/length of vector is 250), each field's value referring to its respective feature. Therefore, with high dimensionality, there is a need to have a large amount of data per vector, as such with dimension-increase, the performance required for accurate results by ML algorithms requires an exponential-increase in data points, since for a feature which can have 2 possible values (binary), assuming 20 data points are needed for the model to perform accurately, and there is 1 dimension, then the number of unique values is $2^1$, then applying this to the number of data points needed (20), $2^1\cdot20 = 40$, in general $numValsPerFeature^{numDimensions}\cdot numDataPointsNeeded$ (in this example, $2^k\cdot 20$, where $k$ is the number of dimensions for these parameters (values per feature, and minimum data points required)). Therefore, adding another dimension results in $2^2\cdot20 = 80$, and once more with 3 dimensions $2^3\cdot20 = 160$, this showing the exponential increase (in this example) in the data required for this model to make accurate predictions \citep{towardsdatascience:curse_of_dimensionality}. %todo if really do have time, then talk about hughes phenomenon here, and perhaps draw graph; refer to citation towardsdatascience:curse_of_dimensionality
\paragraph{Reducing Dimensionality via PCA}%TODO add more detail at end of para.
\label{para:reducing_dimensionality_via_pca}
\noindent\\
Given some data containing a large number of features, PCA transforms these into new features called principal (leading; ``most important'') components, which are based on direction (eigenvector) and perceived importance (eigenvalues) of the features \citep{geeks:pca}, calculating an eigenvectors and corresponding eigenvalues from a covariance matrix \textcolor{red}{(comparing how each feature varies with every other feature -- but also including itself (diagonal elements in covariance matrix showing variance within the given feature itself -- deviation of its points from mean point))}.\\
Now the principal components have been calculated (eigenvectors of covariance matrix), each with importance (eigenvalues), find the eigenvector with the highest eigenvalue (choosing the direction (eigenvector) with the greatest amount of variance, since this indicates useful information, thereby capturing this when projected onto the original data). Projecting the eigenvectors over the original data (visualisation) shows the principal components, taking away the original data -- reduced dimensionality to the chosen, most-relevant, principal components -- thereby retaining the deemed most-useful information, whilst discarding all else. The more principal components that are chosen to be retained, the more the reduced data will resemble the original data (i.e. higher dimensionality, but more data) \citep{geeks:pca}.
\subsection{Backpropagation}
\label{subsec:backpropagation}
Backpropagation is used in training neural networks (NNs) to minimise the \textcolor{red}{cost function -- the difference between overall model output and expected output}, \textcolor{red}{for some given training input data}, this is calculated via a metric showing the error, such as the Mean Squared Error (MSE) (see \cref{eqtn:mse}), which measures the average squared difference between the actual value (in the training data) and the predicted value (from the network) \citep{geeks:mse} -- a lower MSE therefore shows the network's prediction is closer to the actual/expected value; greater prediction accuracy of the model \citep{geeks:mse}. MSE is useful as a cost function, given it is influenced by large errors, more so than small errors (due to squaring the error), focusing on \textcolor{red}{the largest error produced by the network, rather than small errors}. Calculating the gradients of each layer, backwards (from output layer to input layer), uses the chain rule to calculate partial derivatives of each layer, with respect to each layer before (closer to the output layer of the network), this enables determining how much an adjustment in a given layer's weights will impact the network's overall output, and, by extension, how that impact the cost function.\\
This occurs by adjusting the weights of neurons by calculating the gradient per layer backwards (from output layer through the hidden layers, to the input layer) through the network; backward pass, see below. The amount to adjust each layer's weights is determined by the calculated gradient, with respect to each previous layer (where previous layers are those closer to the final layer of the network, given this is the \textbf{backwards} pass) \citep{analyticsvidhya:gradient-descent_backpropagation}.
\begin{center}
    \begin{equation}\label{eqtn:mse}
        \frac{1}{n}\sum_{i=1}^{n}(Y_i-\hat{Y}_{i})^{2}
    \end{equation}
    \vspace{3mm}
    Where $n$ is the number of data points in the training set\footnote{As you may have spotted, with large training datasets, $n$ will be very large, if running this $n$ times per epoch, this is very computationally expensive. One optimisation, therefore, is stochastic  gradient descent, which uses only a (pseudo-)randomly- chosen subset of training data points, $n$ is then the size of this subset of training data.}, $Y_i$ is the actual value of the $i$th data point, and $\hat{Y}_i$ is the predicted value of the $i$th data point \citep{geeks:mse}.
\end{center}
\subsubsection{Gradient Descent}
Given the gradients calculated from backpropagation, gradient descent determines the amount to adjust the weights of a given layer in an attempt to reduce the cost function (until the point at which this has been minimised, \textcolor{red}{i.e. convergence has been reached}).\\
The (inverse of the) gradient calculated by backpropagation determines the ``direction to move in'', i.e. to increase or decrease a given weight -- how much to adjust the weight is determined by a hyperparameter known as the learning rate. If the learning rate value is high, the step size will be large, which presents issues when the network is close to convergence, given the weights may be adjusted by the amount of the learning rate too much so that the output ``bounces'' around the global minimum; if the learning rate is too low, the learning process will take a considerable amount of time (per epoch (/iteration of training, \textcolor{red}{i.e. per both forwards and backwards pass}), little change is made; higher number of epochs is, hence, required).
\citep{analyticsvidhya:gradient-descent_backpropagation}

%gradient-calculation: (partial) derivative of cost function "with respect to each parameter"(chain rule). recalc'd at each layer (presumably during backward pass), the gradient calc'd at each layer shows the (projected) change to the overall loss function given some weight(/bias?)-change

%GD. uses calc'd. gradient to determine the adjustments to weighting (and bias; explain bias!). iteratively done dep. on learning rate (/step size), to cause weighting-movement in the direction+scalar that minimises costFunction.


%then \textcolor{red}{using an algorithm like gradient descent, or one of its derivatives, such as \textcolor{red}{the more- computationally-fast stochastic- gradient-descent}},

%this uses the X calculated to determine how much a change in the weight of a given neuron in a given layer will affect the overall model-output; the weights are thence adjusted to minimise the cost function for the overall model-output


%TODO check accuracy+relevance: so if running BP, and cost-function high (ie output-layer loss is very high) and is not changing (very much at all) despite backpropagating the changes thoruygh the network, well little change would suggest convergence, but if output-loss still high, likely vanishing-gradient problem ig? \citep{geeks:backpropagation}